\documentclass[11pt]{article}
\usepackage[a4paper,margin=2.5cm]{geometry}
\usepackage{amsmath,amssymb}
\usepackage{graphicx}
\usepackage{booktabs}
\usepackage{hyperref}
\usepackage{parskip}

\title{MA-INF 2201 Computer Vision WS26/27 \\ Exercise Sheet 1 --- Written Answers}
\author{Mohamadhosein (Parsa) Sabaghi}
\date{\today}

\begin{document}
\maketitle

% ==========================================================================
\section{Denoising and Optimization}

\subsection{Discussion (Part 1.4)}
% TODO: adjust this once you've visually compared the Part 1.1 / 1.3 result
% figures (see the last sentence below).

After parameter optimization, the Gaussian (22.94~dB) and bilateral
(22.86~dB) filters performed similarly and clearly outperformed the median
filter (22.06~dB) on PSNR. This is somewhat counter to the naive expectation
that median filtering, being ideal for salt-and-pepper noise, should
dominate; however, median filters are statistically less efficient than
averaging-based filters at suppressing the additive Gaussian component of
the noise (the sample median has higher variance than the sample mean for
Gaussian-distributed data), so on this mixed-noise image the Gaussian
component's contribution to the error outweighs median's advantage on the
impulse component. The near-tie between Gaussian and bilateral further
suggests that PSNR, being a purely pixel-wise error metric, does not
strongly reward bilateral's edge-preservation property --- it favors
whichever filter reduces variance most uniformly, regardless of whether
real edges are blurred along with the noise. Visually, we would still
expect bilateral to preserve edges better than Gaussian despite the similar
PSNR, since PSNR does not fully capture perceptual/structural quality.
% TODO: keep or remove the sentence above based on what you actually see in
% the displayed images.

% ==========================================================================
\section{Convolution of Gaussians}

\subsection{Convolution of Two 1D Gaussians (Part 2.1)}

By definition of convolution,
\[
(g_{\sigma_1} * g_{\sigma_2})(x) = \int_{-\infty}^{\infty} g_{\sigma_1}(u)\, g_{\sigma_2}(x-u)\, du
= \frac{1}{2\pi\sigma_1\sigma_2} \int_{-\infty}^{\infty} \exp\!\left(-\frac{u^2}{2\sigma_1^2} - \frac{(x-u)^2}{2\sigma_2^2}\right) du.
\]

Write $A = \dfrac{1}{2\sigma_1^2}$, $B = \dfrac{1}{2\sigma_2^2}$, so the exponent is
$-Au^2 - B(x-u)^2 = -(A+B)u^2 + 2Bxu - Bx^2$.
Completing the square in $u$:
\[
-(A+B)u^2 + 2Bxu - Bx^2 = -(A+B)\left(u - \frac{Bx}{A+B}\right)^2 - \frac{AB}{A+B}x^2.
\]

Substituting back and using the hint $\int_{-\infty}^{\infty} e^{-a(u-b)^2}\,du = \sqrt{\pi/a}$ with $a = A+B$:
\[
(g_{\sigma_1} * g_{\sigma_2})(x) = \frac{1}{2\pi\sigma_1\sigma_2}\, \sqrt{\frac{\pi}{A+B}}\; \exp\!\left(-\frac{AB}{A+B}x^2\right).
\]

Now simplify the two factors. Since $A+B = \dfrac{\sigma_1^2+\sigma_2^2}{2\sigma_1^2\sigma_2^2}$, writing $\sigma^2 := \sigma_1^2+\sigma_2^2$ gives
\[
\frac{AB}{A+B} = \frac{1}{2(\sigma_1^2+\sigma_2^2)} = \frac{1}{2\sigma^2},
\qquad
\sqrt{\frac{\pi}{A+B}} = \sqrt{2\pi}\,\frac{\sigma_1\sigma_2}{\sigma}.
\]

Substituting both back:
\[
(g_{\sigma_1} * g_{\sigma_2})(x) = \frac{1}{2\pi\sigma_1\sigma_2}\cdot \sqrt{2\pi}\,\frac{\sigma_1\sigma_2}{\sigma}\, \exp\!\left(-\frac{x^2}{2\sigma^2}\right)
= \frac{1}{\sqrt{2\pi}\,\sigma}\exp\!\left(-\frac{x^2}{2\sigma^2}\right) = g_\sigma(x),
\]
with $\sigma = \sqrt{\sigma_1^2+\sigma_2^2}$, which is exactly $g_\sigma$. $\blacksquare$

\textbf{Conclusion.} Setting $\sigma_1 = \sigma_2 = \sigma$ gives $(g_\sigma * g_\sigma)(x) = g_{\sigma\sqrt{2}}(x)$: convolving twice with a Gaussian of standard deviation $\sigma$ is identical to convolving once with a Gaussian of standard deviation $\sigma\sqrt{2}$, since both produce a Gaussian whose variance is the sum of the variances of the two convolved kernels ($\sigma^2+\sigma^2 = 2\sigma^2$).

\subsection{Repeated Small Filter vs.\ One Large Filter (Part 2.2)}

\textbf{Number of repetitions.} By the composition rule from Part 2.1, applying a Gaussian of standard deviation $\sigma_1$ $n$ times in a row yields an effective standard deviation $\sigma_{\text{eff}} = \sigma_1\sqrt{n}$ (variances add). With $\sigma_1 = 1$ and the target $\sigma_{\text{eff}} = 4$:
\[
4 = 1\cdot\sqrt{n} \;\Longrightarrow\; n = 16.
\]
So the $\sigma=1$ filter must be applied \textbf{16 times}.

\textbf{Kernel sizes.} Using the rule of thumb half-width $\approx 3\sigma$, i.e.\ full kernel width $w = 6\sigma+1$:
\[
w_{\text{small}} = 6(1)+1 = 7, \qquad w_{\text{large}} = 6(4)+1 = 25.
\]

\textbf{Cost comparison (multiplications per pixel).}

\begin{table}[h]
\centering
\begin{tabular}{lcc}
\toprule
 & Full 2D kernel & Separable (1D $\times$ 1D) \\
\midrule
Repeated small filter ($n=16$ applications of $w=7$) & $16 \cdot 7^2 = 784$ & $16 \cdot (2\cdot 7) = 224$ \\
Single large filter ($w=25$) & $25^2 = 625$ & $2\cdot 25 = 50$ \\
\bottomrule
\end{tabular}
\end{table}

\textbf{Which is cheaper?} In both cases the single large filter is cheaper, but the gap is far more dramatic for the separable implementation ($50$ vs.\ $224$, a $\sim4.5\times$ difference) than for the full 2D kernel ($625$ vs.\ $784$, only $\sim1.25\times$). This is because for a full 2D kernel, repeating a small filter $n=(\sigma_{\text{target}}/\sigma_1)^2$ times costs $n\cdot w_{\text{small}}^2 \approx 36\,\sigma_{\text{target}}^2$, the same asymptotic order as a single large kernel's $w_{\text{large}}^2\approx 36\,\sigma_{\text{target}}^2$ cost --- so the two scale identically with $\sigma_{\text{target}}$, and only the rounding/$+1$ terms separate them. For a \emph{separable} kernel, however, a single pass costs only $2w\approx 12\,\sigma_{\text{target}}$ (linear in $\sigma_{\text{target}}$), while repeating the small separable filter still costs $n\cdot 2w_{\text{small}} \approx 12\,\sigma_{\text{target}}^2/\sigma_1$ (quadratic in $\sigma_{\text{target}}$, since $n$ itself grows quadratically). So whenever separable filtering is available, using one correctly-sized large filter is substantially cheaper than cascading many small ones --- the advantage of separability compounds with the size of the single large filter, but is largely wasted by repeating a tiny filter many times.

% ==========================================================================
% TODO: add Section 3 (Integral Images) back once you're ready.
% TODO: add Section 4 (Histogram Equalization) back once you're ready.

\end{document}
